Phần này liệt kê và giới thiệu các bài blog mình đã đăng hoặc đang chuẩn
bị đăng trong nhóm
\href{https://www.facebook.com/groups/awsstudygroupfcj}{AWS Study
Group}. Các bài blog được chia thành ba mục nhỏ, giữ cùng format với
trang workshop mẫu.

\subsubsection{\texorpdfstring{\href{3.1-Blog1/}{Blog 1 - TẮT EC2 RỒI
NHƯNG AWS VẪN TÍNH
PHÍ?}}{Blog 1 - TẮT EC2 RỒI NHƯNG AWS VẪN TÍNH PHÍ?}}

Bài blog này giải thích lý do AWS vẫn có thể phát sinh chi phí dù EC2
instance đã được tắt, tập trung vào các tài nguyên liên quan vẫn có thể
tiếp tục được tính phí ngoài bản thân instance đang chạy.

\textbf{Ngày đăng:} July 22, 2026\\
\textbf{Nền tảng:} AWS Study Groups\\
\textbf{Trạng thái:} Approval\\
\textbf{Public link:}
\href{https://www.facebook.com/groups/awsstudygroupfcj/permalink/2219932025438424/?rdid=SGRgenG5JS3UR6Y0\#}{Blog
1}

\subsubsection{\texorpdfstring{\href{3.2-Blog2/}{Blog 2 - XÂY DỰNG HỆ
THỐNG CHAT REALTIME CÓ KHẢ NĂNG MỞ RỘNG VỚI SOCKET.IO, REDIS VÀ
DYNAMODB}}{Blog 2 - XÂY DỰNG HỆ THỐNG CHAT REALTIME CÓ KHẢ NĂNG MỞ RỘNG VỚI SOCKET.IO, REDIS VÀ DYNAMODB}}

Bài blog này giới thiệu cách thiết kế một hệ thống chat realtime có khả
năng mở rộng bằng Socket.IO, Redis và DynamoDB. Nội dung tập trung vào
cách giao tiếp thời gian thực, trạng thái phân tán và lưu trữ lâu dài có
thể phối hợp trong một hệ thống lớn hơn.

\textbf{Ngày đăng:} July 23, 2026\\
\textbf{Nền tảng:} AWS Study Groups\\
\textbf{Trạng thái:} Pending\\
\textbf{Public link:} Blog 2

\subsubsection{\texorpdfstring{\href{3.3-Blog3/}{Blog 3 -
\ldots{}}}{Blog 3 - \ldots{}}}

Bài blog thứ ba chưa được hoàn thành. Nội dung sẽ được cập nhật sau khi
chủ đề được chốt và bài viết được đăng.
